\chapter{The Lamphrodynamic Plenum: Axiomatic Foundation}

\section{Motivation and Overview}

The purpose of this chapter is to establish a rigorous axiomatic framework
for the lamphrodynamic plenum, consisting of a triple
$(\Phi,\vec v,S)$ of scalar, vector, and entropy fields.
The physical interpretation of these fields has been discussed informally
in earlier writings; here we develop a formal foundation suitable for
derived geometry and subsequent quantization via the AKSZ construction.

The lamphrodynamic plenum is conceived as a universal medium
of information and energetic exchange.
The scalar potential $\Phi$ encodes lamphrodynamic potential energy,
the vector field $\vec v$ represents oriented transport flows, and
the entropy field $S$ accounts for dissipative and diffusive structure.
The fundamental equations couple these fields in a manner that respects
derived symplectic geometry, thermodynamic irreversibility, and
causal structure.

The classical field theory of $(\Phi,\vec v,S)$ will be defined by a
variational principle in Section~5.5, producing Euler--Lagrange
equations whose analytical properties will be investigated in
Chapters~9--11. Our goal in this chapter is to introduce the field
variables, state the axioms that govern their behavior, and derive
initial differential identities and structural constraints.


\section{Definition of the Lamphrodynamic Fields}

Let $(M,g)$ be a smooth oriented time-oriented Lorentzian manifold
of dimension $4$. The lamphrodynamic fields consist of:

\begin{enumerate}
\item A real-valued scalar field $\Phi \in C^{\infty}(M,\mathbb{R})$.
\item A spacetime vector field $\vec v \in \Gamma(TM)$.
\item A real-valued entropy scalar field $S \in C^{\infty}(M,\mathbb{R})$.
\end{enumerate}

We write collectively
\[
\mathcal{U} := (\Phi,\vec v,S).
\]

\begin{axiom}[Field regularity]
For the purposes of the classical theory developed in this volume,
we assume $\Phi$, $S$ are smooth and $\vec v$ is a smooth vector field
on $M$. Weak and distributional extensions will be introduced in
Chapter~11.
\end{axiom}

\begin{remark}
In certain applications we will regard $S$ as taking values in
$\mathbb{R}_{\ge 0}$, though for the mathematical theory it suffices to
treat $S$ as real-valued.
\end{remark}


\section{Lamphrodynamic Coupling Principles}

The physical behavior of the scalar, vector, and entropy fields derives
from the principle that lamphrodynamics is governed by coupled 
transport, diffusion, and potential flow. We impose the following axioms.

\begin{axiom}[Transport--Potential coupling]
The vector field $\vec v$ transports the scalar potential $\Phi$
according to the advection equation
\begin{equation}
\label{eq:advect}
\mathcal{L}_{\vec v}\Phi = -\kappa_{\Phi}\Delta_g \Phi + R_{\Phi},
\end{equation}
where $\kappa_{\Phi}\ge 0$ is a diffusion coefficient,
$\Delta_g$ is the Laplace--Beltrami operator, and $R_{\Phi}$
is a nonlinear lamphrodynamic coupling term defined in
Section~5.6.
\end{axiom}

\begin{axiom}[Entropy production]
The entropy field $S$ satisfies
\begin{equation}
\label{eq:entropy}
\mathcal{L}_{\vec v}S = \kappa_{S}\Delta_g S + \sigma(\Phi,\vec v,S),
\end{equation}
where $\sigma(\Phi,\vec v,S)\ge 0$ is the entropy production rate.
\end{axiom}

\begin{axiom}[Velocity response]
The transport field $\vec v$ responds to both entropy and potential
gradients by
\begin{equation}
\label{eq:velocity}
\nabla_{\vec v}\vec v = -\nabla\Phi + \beta\nabla S + R_{v},
\end{equation}
where $\beta\in\mathbb{R}$ couples entropy gradients to lamphrodynamic
flows and $R_{v}$ is described in Section~5.7.
\end{axiom}

These axioms are formal for now; precise variational origins will be
developed in Section~5.5. The choice of signs corresponds to the
interpretation of $-\nabla\Phi$ as attractive potential flow and
$\beta\nabla S$ as repulsive entropic response.


\section{Lamphrodynamic Conservation and Balance Laws}

We now derive several differential identities and balance laws implied
by the axioms of the previous section. For clarity, we work in local
coordinates with the Levi--Civita connection of the Lorentzian metric $g$.

\begin{proposition}[Entropy balance]
Assume equation \eqref{eq:entropy}. Then
\begin{equation}
\partial_t \int_{\Sigma_t} S \, d\mu_t
= \int_{\Sigma_t} \left( \kappa_{S}\Delta_g S
  + \sigma(\Phi,\vec v,S) \right) d\mu_t
\end{equation}
for any spacelike hypersurface $\Sigma_t$ with induced measure $d\mu_t$.
\end{proposition}

\begin{proof}[Proof]
Apply the Reynolds transport theorem to \eqref{eq:entropy} and use
$\mathcal{L}_{\vec v}S = \vec v(S)$.
\end{proof}

Similarly, we have a potential energy balance:

\begin{proposition}[Potential balance]
If $\Phi$ satisfies \eqref{eq:advect}, then
\begin{equation}
\partial_t \int_{\Sigma_t} \Phi \, d\mu_t
= \int_{\Sigma_t}
\left( -\kappa_{\Phi}\Delta_g \Phi + R_{\Phi} \right) d\mu_t.
\end{equation}
\end{proposition}

The vector field $\vec v$ obeys a momentum-like balance obtained by
integrating \eqref{eq:velocity} against $\vec v$ and using the identity
$\frac12\nabla_{\vec v}\|\vec v\|^2 = \langle\nabla_{\vec v}\vec v,\vec v\rangle$.

\begin{proposition}[Lamphrodynamic momentum balance]
Assuming \eqref{eq:velocity},
\begin{equation}
\frac12 \mathcal{L}_{\vec v}\|\vec v\|^2 =
-\langle\nabla\Phi,\vec v\rangle
+ \beta\langle\nabla S,\vec v\rangle
+ \langle R_{v},\vec v\rangle.
\end{equation}
\end{proposition}

\begin{remark}
The quantity $\langle\nabla S,\vec v\rangle$ plays the role of
entropic tension, corresponding physically to emergent repulsion
induced by entropy gradients. This term will be crucial in the
cosmological applications of Volume~II.
\end{remark}

\section{Variational Principle and Euler--Lagrange Equations}
\label{sec:variational}

We now give the variational origin of the lamphrodynamic equations.
Let $\mathcal{L}$ be a Lagrangian density on $(M,g)$ of the form
\begin{equation}
\label{eq:Lagrangian}
\mathcal{L}(\Phi,\vec v,S;g)
= \tfrac12\|\vec v\|^2
+ U(\Phi) + W(S) + \alpha\langle\nabla S,\vec v\rangle
+ \beta\langle\nabla\Phi,\vec v\rangle,
\end{equation}
where $U$ and $W$ are smooth potential functions and
$\alpha,\beta\in\mathbb{R}$ are coupling constants. Let
\[
\mathcal{S}[\Phi,\vec v,S]
:= \int_M \mathcal{L}\, dV_g
\]
be the associated action functional.
Variation against $\Phi$, $S$, and $\vec v$ produces
Euler--Lagrange equations which we now compute formally.

\vspace{0.5em}
\noindent{\bf Variation with respect to $\Phi$.}
Integrating by parts yields
\[
\delta_{\Phi}\mathcal{S}
= \int_M \left(
U'(\Phi) + \beta \operatorname{div}(\vec v)
\right)\delta\Phi\, dV_g,
\]
and hence the Euler--Lagrange equation is
\begin{equation}
\label{eq:EL-Phi}
U'(\Phi) + \beta\operatorname{div}(\vec v) = 0.
\end{equation}

\vspace{0.5em}
\noindent{\bf Variation with respect to $S$.}
Similarly,
\[
\delta_{S}\mathcal{S}
= \int_M \left(
W'(S) + \alpha \operatorname{div}(\vec v)
\right)\delta S\, dV_g,
\]
producing
\begin{equation}
\label{eq:EL-S}
W'(S) + \alpha\operatorname{div}(\vec v)=0.
\end{equation}

\vspace{0.5em}
\noindent{\bf Variation with respect to $\vec v$.}
A computation using the Levi--Civita connection gives
\[
\delta_{\vec v}\mathcal{S}
= \int_M
\left(
\vec v + \alpha\nabla S + \beta\nabla\Phi
\right)\cdot \delta\vec v\, dV_g,
\]
leading to
\begin{equation}
\label{eq:EL-v}
\vec v + \alpha\nabla S + \beta\nabla\Phi = 0.
\end{equation}

\begin{remark}
Equations \eqref{eq:EL-Phi}, \eqref{eq:EL-S}, and \eqref{eq:EL-v}
represent the stationary points of the lamphrodynamic action.
Their consistency with the axioms of Section~5.3 requires suitable
choices of $U,W,\alpha,\beta$; Section~5.6 analyzes this structure.
\end{remark}


\section{Nonlinear Couplings and Constraint Relations}
\label{sec:couplings}

We now discuss the nonlinear terms $R_{\Phi}$ and $R_{v}$ introduced
formally in Section~5.3. These terms encode constraint relations among
$\Phi$, $\vec v$, and $S$ that are not captured by the minimal action
\eqref{eq:Lagrangian}. In physical terms, $R_{\Phi}$ and $R_{v}$
represent lamphrodynamic self-interaction and dissipation beyond the
lowest-order coupling.

\begin{axiom}[Lamphrodynamic constraint]
There exist smooth functions $\mathcal{C}_\Phi(\Phi,\vec v,S)$ and
$\mathcal{C}_v(\Phi,\vec v,S)$ such that
\begin{equation}
\label{eq:R-Phi}
R_{\Phi} = \mathcal{C}_\Phi(\Phi,\vec v,S),\qquad
R_{v} = \mathcal{C}_{v}(\Phi,\vec v,S),
\end{equation}
subject to the compatibility condition
\begin{equation}
\label{eq:compatibility}
\operatorname{div}(R_{v})
+ \langle\nabla S,R_{v}\rangle
+ \langle\nabla\Phi,R_{v}\rangle
= \Delta_g R_{\Phi}.
\end{equation}
\end{axiom}

\begin{remark}
Condition \eqref{eq:compatibility} guarantees that
Lamphrodynamic balance laws are preserved at nonlinear order,
and plays a role analogous to constitutive relations in hydrodynamics.
\end{remark}


\section{Symmetries and Noether Currents}
\label{sec:symmetry-noether}

We discuss the symmetries of the lamphrodynamic plenum at the level
of the classical action \eqref{eq:Lagrangian}.

\begin{definition}
A local symmetry of the lamphrodynamic action is a variation
$(\delta\Phi,\delta\vec v,\delta S)$ such that
$\delta\mathcal{S}=0$ for all field configurations satisfying
appropriate boundary conditions.
\end{definition}

\begin{proposition}[Shift symmetry of $\Phi$]
If $U$ depends only on $\nabla\Phi$ (e.g.\ pure gradient theory),
then the transformation $\Phi\mapsto \Phi+\epsilon$ is a symmetry
of $\mathcal{S}$, yielding a conserved current
\[
J^{\mu}_{\Phi} = \frac{\partial\mathcal{L}}{\partial(\nabla_{\mu}\Phi)}.
\]
\end{proposition}

\begin{remark}
In cosmological settings, $U$ will typically depend on $\Phi$ itself,
breaking shift symmetry and producing effective potential-induced flow.
\end{remark}

Likewise, if $W$ depends only on $\nabla S$, then $S\mapsto S+\epsilon$
is a symmetry; otherwise, entropy production breaks $S$-shift symmetry.
The behavior of these symmetries in the presence of $R_{\Phi}$ and
$R_v$ is subtle and will be discussed further in Volume~II.


\section{Conservation Laws and Energy Conditions}

We conclude this chapter with an analysis of energy conditions for
lamphrodynamic flows.
Let
\[
\mathcal{E} := \tfrac12\|\vec v\|^2 + U(\Phi) + W(S)
\]
denote the lamphrodynamic energy density.
Assuming equations \eqref{eq:EL-Phi}--\eqref{eq:EL-v}, one computes

\begin{proposition}[Lamphrodynamic energy balance]
\begin{equation}
\mathcal{L}_{\vec v}\mathcal{E}
= - \alpha\|\nabla S\|^2 - \beta\|\nabla\Phi\|^2
+ \langle R_{v},\vec v\rangle + R_{\Phi}.
\end{equation}
\end{proposition}

\begin{proof}[Proof]
Differentiating $\mathcal{E}$ along the flow of $\vec v$ and substituting
\eqref{eq:EL-Phi}--\eqref{eq:EL-v} yields the stated identity.
\end{proof}

\begin{remark}
If $\alpha,\beta>0$ and $R_{\Phi}=R_{v}=0$, then
$\mathcal{L}_{\vec v}\mathcal{E}\le 0$, indicating lamphrodynamic
dissipation. The nonlinear regime will be addressed in Chapter~10 via
a priori energy estimates.
\end{remark}

\medskip

\noindent\textbf{Summary and Preview.}
This chapter has introduced the lamphrodynamic fields, stated axioms for
their coupling, developed a variational formulation, and identified
initial balance laws and symmetries. In Chapter~6 we construct the
derived moduli stack of lamphrodynamic configurations and show that it
admits a canonical $(3-1)=2$-shifted symplectic structure,
crucial for the AKSZ quantization developed in Chapter~7.

